\chapter{面向AI辅助开发的上下文工程方法}


\section{总体框架}
（本节正文待撰写：三层框架——流程约束层/知识组织层/上下文生成层，及与RQ1--RQ4的问题映射。）


\section{流程约束架构}
（本节正文待撰写：需求->设计->计划->测试->实现全过程约束，BRD->PRD->Explore->Design->Plan->Test->Apply，三段式契约与状态留痕。）


\section{双层知识组织模型}
本章阐述 wpw 系统中项目知识的组织、建模与消费方法，对应本文创新点 2（面向智能体上下文生成的能力-代码双层知识组织方法）与创新点 3（图谱驱动的任务上下文压缩方法）。本章首先介绍知识库的构建，即项目知识的采集、组织与稳态沉淀机制（3.1 节）；其次介绍知识图谱的构建，包括数据来源与多语言解析处理（3.2.1 节）、基于多源证据融合的业务-代码关联标注（3.2.2 节），以及图谱的本体设计、实体与关系抽取、存储与增量更新（3.2.3 节）；最后阐述知识库和知识图谱与大语言模型的结合方法，即面向智能体任务上下文的检索、裁剪与压缩流水线（3.3 节）。

需要预先说明本章的知识建模定位：与以代码本身为建模对象的代码属性图（Code Property Graph）[18]等代码结构图谱不同，wpw 知识图谱在代码结构层之上引入了业务能力层，将业务能力节点作为上层语义锚点，使图谱不仅能够回答"代码结构是什么"，还能够回答"这段代码服务于什么业务能力"。这一定位源于第 1 章的研究空白：智能体任务上下文的生成需要业务语义与代码结构的双重知识，而非单纯的代码关系图。


\subsection{知识库的构建}
wpw 的知识基础设施自底向上分为三层：数据来源层（能力规范、需求过程文档、源码、Git 历史）、知识组织层（稳态知识库与知识图谱）与智能体消费层（任务上下文生成），整体架构如图 3-3 所示。数据来源经知识处理流水线（解析、结构化、证据融合）沉淀为知识组织层的稳态知识，再经上下文生成流水线（3.3 节）转化为智能体可消费的任务上下文。3.1 节阐述前两层的构建方法。

\begin{quote}
\textbf{图 3-3 wpw 知识库整体架构}（成稿时绘制）：能力规范 -> 稳态知识库；源码/需求文档/Git 历史 -> 知识处理流水线 -> 知识图谱 -> Agent Context 三层结构。
\end{quote}


\subsubsection{数据采集}
wpw 将项目知识按其生命周期特征分为稳态知识与过程知识两类，并从四个渠道采集。

\textbf{（1）能力规范采集（稳态知识）}。以 OpenSpec 规格格式组织的能力规范存放于 \texttt{wpw/specs/} 目录下，每个能力对应一个子目录及其 \texttt{spec.md} 文件，采用 Purpose（能力目的）、Requirements（需求条目）、Scenarios（验收场景）三段式结构。能力规范描述项目长期存在的业务能力（如"用户认证""订单管理"），是知识图谱 C 层的唯一数据来源，也是业务-代码关联映射的语义基准。

\textbf{（2）需求过程文档采集（过程知识）}。六阶段工作流（BRD->PRD->Design->Plan->Test->Apply）在执行过程中产出需求文档、能力规格提案、技术方案、开发计划与测试方案等阶段性文档，存放于 \texttt{wpw/active/<需求名>/} 目录。过程知识服务于当前开发任务，在需求归档后经沉淀机制转化为稳态知识。

\textbf{（3）源码采集（结构知识）}。项目源码文件是知识图谱结构层（L1/L2/L3）的数据来源。wpw 支持多语言源码的采集与解析，包括 TypeScript/TSX、JavaScript/JSX、Vue 单文件组件、Java（Spring Boot），并覆盖 Pinia、Vuex、Redux 等前端状态管理生态与微信小程序、uni-app 等平台扩展。

\textbf{（4）Git 历史采集（演化知识）}。项目的提交历史（commit message 与文件变更记录）反映了业务概念与代码文件之间的长期关联强度，作为业务-代码映射的证据来源之一（见 3.2.2 节）。


\subsubsection{知识组织与稳态沉淀}
知识库的组织核心是"过程知识向稳态知识"的沉淀机制，其目标是解决 AI 辅助开发中业务知识随需求结束而流失的问题。

\textbf{（1）能力规范结构化}。对每个能力规范文件，系统从 Purpose 章节提取能力描述（description），从 Requirements 章节结构化提取功能条目数组（features），每个条目包含编号、名称、优先级与描述字段。结构化后的能力规范同时服务于图谱 C 层节点生成（3.2.3 节）与语义向量构建（3.3 节）。

\textbf{（2）稳态能力沉淀机制}。在需求归档（archive）时，系统对已交付需求的业务内容进行能力归属分析：若需求归属于已有能力，则将其增量内容以 delta 方式合并入对应能力规范；若属于新能力领域，则创建新的能力规范文件。该机制使业务知识随项目推进持续积累，且始终以机器可解析的规格形式存在，而非散落在一次性过程文档中。需要说明的是，能力归属分析与合并由 AI 层在流程约束下执行并经用户确认，系统提供目录约定、格式校验与图谱增量更新保障其落盘一致性。

\textbf{（3）存储方案}。知识库采用纯文件系统存储，无数据库依赖：能力规范以 Markdown 加结构化字段存于 \texttt{wpw/specs/}；知识图谱以 JSONL 逐行存储节点与边；语义向量以二进制文件存储并配合映射索引；工作流状态以 \texttt{.wpw.yaml} 存储。文件系统存储使知识库天然具备版本可控（可随代码一同提交至 Git）、人机可读、零部署成本的特性。


\subsection{知识图谱的构建}

\subsubsection{数据来源与处理}
知识图谱的数据来源与分层模型一一对应：能力规范生成 C 层（业务能力）节点；源码目录结构经模块推断生成 L1（模块）节点；源码文件解析生成 L2（文件）节点；文件内的代码元素解析生成 L3（函数、类、接口、常量等）节点。

\textbf{（1）多语言源码解析}。结构层节点的抽取基于 web-tree-sitter 的语法树解析完成。选择 tree-sitter 而非 Babel Parser、TypeScript Compiler API 等单语言 AST 工具作为统一解析基础，基于四点考虑：其一，多语言统一接口--tree-sitter 为不同语言提供一致的节点遍历与字段访问 API，使解析器可按统一模式扩展至 TS/JS/Vue/Java 等语言，而单语言工具需为每门语言维护独立适配层；其二，增量解析能力--tree-sitter 原生支持基于旧语法树的增量重解析，与图谱的增量更新机制（3.2.3 节）天然契合；其三，容错性--其在语法错误处仍能产出局部可用的语法树，适合解析真实项目中不完整或含方言的代码；其四，跨语言一致性--wpw 面向的前后端全栈项目需要统一的多语言解析管线。

在解析策略上，考虑到不同语言的解析成本差异，系统采用"快速预检 + 完整解析"的两级策略：先以轻量级特征检测（路径特征与内容特征，如 Pinia store 文件的 \texttt{defineStore} 调用特征）判断文件是否属于特定扩展类型，命中后再调用对应解析器执行完整的语法树解析，避免对所有文件无差别执行重解析。多语言解析的覆盖范围为：TypeScript/TSX/JavaScript/JSX 提取函数、类、接口、常量与导入关系；Vue 单文件组件提取 script 块并复用 TS/JS 解析；Java（Spring Boot）提取类型（class/interface/enum/record）、方法、常量、注解与 REST 端点元信息；Pinia/Vuex/Redux 提取 store、action、getter、state 等状态管理元素；小程序与 uni-app 提取页面、组件、生命周期与路由关系。

\textbf{（2）节点上下文属性与去歧义}。同名代码元素（如多个组件中的 \texttt{onSubmit} 函数）是代码检索与语义匹配的主要噪声来源。为此，所有 L3 元素节点在生成时强制携带两类上下文属性：\texttt{filePath}（所属文件相对路径）与 \texttt{parentName}（所属上层容器名，如类名、store 名或模块文件名）。该设计使后续的向量富化（3.3 节）与检索结果展示能够以"容器/元素"的形式区分同名实体，其有效性在消融实验的图谱贡献组（5.1.5 节消融 A）中验证。L3 节点的完整属性定义如表 3-3 所示，后续的向量构建、语义检索与上下文序列化均依赖该属性模式。

表 3-3 L3 代码元素节点属性定义

\begin{table}[htbp]
\centering
\small
\begin{tabularx}{\textwidth}{|X|X|X|}
\hline
\textbf{字段} & \textbf{类型} & \textbf{说明} \\
\hline
id & string & 节点唯一标识（类型前缀 + 内容哈希） \\
type & enum & function / class / interface / constant 及状态管理元素类型 \\
name & string & 元素名称（方法节点带父容器前缀，如 \texttt{UserService.findById}） \\
filePath & string & 所属文件相对路径（去歧义） \\
parentName & string & 父容器名：类名/store 名/模块文件名（去歧义） \\
signature & string & 函数签名或类型声明摘要 \\
params & array & 参数列表（名称与类型） \\
returnType & string & 返回类型 \\
annotations & array & 装饰器/注解（Java Spring 注解等） \\
jsDoc & string & JSDoc/Javadoc 注释摘要（参与向量富化） \\
lineStart/lineEnd & number & 源码行号范围 \\
\hline
\end{tabularx}
\end{table}


\subsubsection{数据标注：多源证据融合的业务-代码映射}
知识图谱构建的关键问题是建立 C 层业务能力与 L1/L2/L3 结构层节点之间的关联（business\_map 边）。需求-代码追踪研究表明，人工建立并维护此类关联的成本高昂[13-14]，自动追踪方法则长期面临准确率与召回率的权衡[15]。wpw 的定位并非提出新的追踪算法，而是面向智能体上下文生成场景，将多类启发式证据以概率模型融合，以较低成本获得"可用且可溯源"的业务-代码关联，并通过权重剪枝控制噪声。business\_map 边的生成流程为：能力节点 -> 候选召回 -> 四类证据计算 -> noisy-OR 融合 -> 阈值过滤 -> business\_map 边，如图 3-4 所示。

\begin{quote}
\textbf{图 3-4 business\_map 生成流程}（成稿时绘制）：能力节点 -> 候选召回 -> 四类证据并行计算 -> noisy-OR 聚合 -> 权威性溯源标注 -> 阈值过滤 -> 边生成。
\end{quote}

\textbf{（1）证据源设计}。系统实现了四类自动化证据源，各证据源的机制与权重参数如表 3-1 所示。

表 3-1 多源证据体系

\begin{table}[htbp]
\centering
\small
\begin{tabularx}{\textwidth}{|X|X|X|X|}
\hline
\textbf{证据源} & \textbf{机制} & \textbf{基础权重} & \textbf{权重上限} \\
\hline
文档提取（doc-extract） & 从能力规范的 Purpose/Requirements 章节抽取模块名与关键词，直接匹配图谱中的 L1/L2 节点 & 0.85 & 0.85 \\
命名匹配（name-match） & 中英文词典（≥30 词条）翻译匹配 + 前缀/包含匹配，命中能力名与节点名 & 0.45（词典）\textasciitilde{}0.60（直接包含） & 0.70 \\
语义匹配（semantic） & 能力节点与 L1/L2 节点向量的余弦相似度线性映射，每能力取 Top-K（K=5）候选 & 相似度线性映射 & 0.70 \\
Git 历史追溯（git-history） & 以能力名与关键词检索提交历史，统计匹配提交修改文件的频次并归一化（频次 ≥2 方可作为证据） & 频次归一化映射 & 0.70 \\
\hline
\end{tabularx}
\end{table}

四类证据源分别从"文档明示""命名约定""语义相似""演化关联"四个独立视角建立关联假设，其证据特征彼此互补：文档提取权威性最高但覆盖受限于规范书写完整度；语义匹配覆盖面广但受向量表示质量影响；Git 历史反映长期实践关联但对新项目（冷启动）失效。需要说明的是，表中基础权重与上限为启发式经验设置的初始值，其取值以第 5 章组件级基线实验的人工抽样标注准确率为参照进行校准与敏感性检验；本文的重点不在优化追踪模型的参数，而在验证多源证据融合对智能体上下文生成效果的贡献（见消融 A 与组件级基线）。

\textbf{（2）noisy-OR 权重聚合}。同一目标节点被多个证据源命中时，采用 noisy-OR 概率模型聚合各源权重：

\begin{equation}
W = 1 - \prod_{i=1}^{n}(1 - w_i) \tag{3-1}
\end{equation}

其中 $w_i$ 为第 $i$ 个证据源对该目标的基础权重，$W$ 为聚合权重，上限截断为 0.95。采用 noisy-OR 的依据是：将各证据源视为对"该节点属于该业务能力"这一事件的独立概率化支持，多源命中的可信度应严格高于任一单源（除非已达上限），该性质与证据互补的设计意图一致。例如，某文件节点同时被文档提取（$w_1=0.85$）与语义匹配（$w_2=0.60$）命中，则聚合权重 $W = 1-(1-0.85)(1-0.60) = 0.94$，高于任一单源。相对于最大权重、加权平均等聚合方式，noisy-OR 对多源确认更敏感且无需预设源间权重超参；三种聚合方式的定量对比见第 5 章消融实验的组件级基线部分。

\textbf{（3）证据权威性排名与溯源}。不同证据源的可信度存在差异，系统为每条 business\_map 边记录支撑该边的最权威证据源（source 字段），权威性排名为：structure（结构证据，rank 10）> doc-extract（8）> ai-refine（7，预留）> git-history（5）> semantic（4）> name-match（2）。该排名在多源命中时决定边溯源标签的取值，使每条业务映射边可回溯至其最可信的证据来源，支撑第 5 章可追溯性维度的评估。

\textbf{（4）低权重剪枝}。聚合权重低于 0.3 的候选关联不生成边，以控制图谱噪声规模。剪枝阈值与各源权重参数共同构成映射质量的调控面，其取值以第 5 章组件级基线实验的抽样人工标注准确率为参照校准。


\subsubsection{知识图谱的构建}
\textbf{（1）本体设计}。wpw 知识图谱采用"能力-结构"双层四层本体模型，如图 3-1 所示：C 层为业务能力节点，表达项目长期稳定的业务语义；L1 层为模块节点（业务模块，带 frontend/backend 侧别属性）；L2 层为文件节点；L3 层为代码元素节点（函数、类、接口、常量及各状态管理生态元素）。

结构层划分为三层而非单层或两层，是因为三个粒度分别对应智能体开发任务中的三类定位需求：模块层承担\textbf{业务边界定位}（影响范围分析--"该需求波及哪些业务模块"），文件层承担\textbf{代码物理定位}（文件检索--"修改发生在哪些文件"），元素层承担\textbf{修改粒度定位}（代码编辑--"具体修改哪个函数/类"）。三层缺一，任务上下文即缺失一级定位能力；C 层则在三者之上提供业务语义锚点，使自然语言任务描述能够经"能力 -> 模块 -> 文件 -> 元素"的路径逐级落地到代码。

以"用户认证"能力为例的图谱实例：C 层节点 \texttt{user-auth}（携带能力描述与功能条目）经 business\_map 边映射至 L1 模块 \texttt{auth}（business\_map 权重 0.85，溯源标签 doc-extract）；\texttt{auth} 模块经 contain 边包含 L2 文件 \texttt{login.vue} 与 \texttt{auth.service.ts}；\texttt{login.vue} 经 contain 边包含 L3 元素 \texttt{handleLogin()}（携带 filePath 与 parentName 属性）；\texttt{auth.service.ts} 中的 \texttt{AuthService.findById} 方法作为 public 方法生成 L3 节点。当智能体接受"修改登录逻辑"类任务时，语义检索命中 \texttt{user-auth} 能力节点后，即可沿该路径展开结构完整的任务上下文。

结构层内部以 contain（包含）边构成层级树（L1 ⊃ L2 ⊃ L3），以 import（导入）边表达文件间依赖，以 call（调用）边表达组件对状态管理 action 的调用；C 层与结构层之间以 business\_map（业务映射）边连接，\textbf{不使用 contain 边}--业务能力与代码结构是"映射"而非"包含"关系，这是本图谱与单一层级代码图谱（如以文件系统包含关系为主体的代码属性图[18]）在语义上的本质区别，也是业务能力节点能够作为上层语义锚点参与检索与裁剪的前提。

表 3-2 图谱边类型

\begin{table}[htbp]
\centering
\small
\begin{tabularx}{\textwidth}{|X|X|X|X|X|}
\hline
\textbf{边类型} & \textbf{语义} & \textbf{方向} & \textbf{权重} & \textbf{来源} \\
\hline
contain & 层级包含（模块⊃文件⊃元素） & 上层 -> 下层 & 0.9 & 结构解析 \\
import & 文件间导入依赖 & 导入方 -> 被导入方 & 0.75 & 结构解析 \\
call & 组件调用状态管理 action & 调用方 -> action & 0.6 & 结构解析 \\
business\_map & 业务能力映射到结构层节点 & C -> L1/L2/L3 & noisy-OR 聚合（≤0.95） & 多源证据融合 \\
\hline
\end{tabularx}
\end{table}

\textbf{（2）实体抽取}。C 层实体由能力规范解析器遍历 \texttt{wpw/specs/} 目录生成，节点名取能力目录名，携带 description 与 features 属性；结构层实体由 3.2.1 节所述多语言解析器从语法树抽取。所有节点 ID 采用"类型前缀 + 内容哈希（SHA-256 前 12 位十六进制）"的确定性生成方式，相同内容必然生成相同 ID，保证图谱重建的幂等性，也为增量更新提供基础。

\textbf{（3）关系抽取}。contain 与 import 关系由结构解析直接导出；business\_map 关系由 3.2.2 节多源证据融合生成。向量语义匹配所需的节点向量在关系抽取前构建：向量输入文本对节点类型分层富化，L3 元素节点的向量文本为"文件路径 + 父容器名 + 节点名 + 注释文本"的拼接，中文注释（JSDoc/Javadoc）的纳入显著提升了中文业务语句与代码元素之间的语义检索效果。

\textbf{（4）图谱存储}。图谱以纯文件系统方案存储于 \texttt{wpw/knowledge/graph/<图谱名>/} 目录：节点与边以 JSONL 逐行存储（graph.jsonl）；语义向量以二进制文件存储（32 字节头记录维度与数量，正文为 Float32Array）并配 JSON 映射索引（向量下标与节点 ID 双向映射）；构建元数据（schema 版本、文件哈希快照、扫描根、项目类型等）存于 meta.json。schema 版本当前为 3.1.0，主版本号不兼容时自动降级为全量重建，次版本升级保持增量兼容。

\textbf{（5）增量更新}。全量重建对中型项目成本较高，系统支持基于文件哈希快照的增量更新：对比当前源码文件哈希与 meta.json 中记录的历史快照，仅对新增与变更文件重新解析并重建其关联节点与出边；同时检测能力规范的增删改，同步更新 C 层节点。已知局限是：被引用方（如被重命名或删除的类）自身的变更不会触发引用它的其他未变更文件的入边重建，可能产生悬挂引用；系统在完整性校验阶段对此输出告警，并建议在结构性重构后执行全量重建。


\section{图谱驱动的任务级上下文生成}

\subsection{知识库和知识图谱与大模型的结合}
\textbf{（1）结合动机：从文本检索到结构感知上下文生成}。常规 RAG 方法将文档切片后以向量相似度检索并拼接注入[5]，该范式应用于代码场景存在两个不适配：其一，代码的语义不仅存在于文本，更存在于结构（调用、包含、导入关系）之中[11,18]，纯文本检索无法保证返回的任务上下文在结构上连贯；其二，检索结果的注入规模缺乏结构化的预算控制，容易重现"Lost in the Middle"[4]的信息利用问题。因此，wpw 的上下文生成方法不采用"检索文档片段"的思路，而是\textbf{在知识图谱上执行结构约束的子图裁剪}：以图结构保证上下文的连通性与层次完整性，以图算法（而非人工整理）控制注入规模。需要强调的是，本方法并非替代向量检索，而是在向量检索提供候选入口（锚点）之后，引入代码结构图约束完成上下文的组织、裁剪与压缩--向量检索解决"从哪里进入图谱"，图结构约束解决"进入后取什么、取多少"。这正是本文创新点 3 与通用 RAG 的核心区别。

上下文生成流水线分为四个阶段：锚点选择 -> 子图裁剪 -> 骨架抽取 -> 序列化输出，如图 3-2 所示，整体流程以算法 3-1 形式化描述。

\textbf{算法 3-1 图谱驱动的任务上下文生成}

\begin{verbatim}
输入：任务描述 q，Token 预算 B，图谱 G = (V, E)
输出：智能体任务上下文 C

1   A ← 语义检索定位锚点候选            // 向量相似度，C 层入口
2   A ← A ∪ business_map 扩展的结构层锚点  // 置信度衰减加权（式 3-3）
3   G' ← ∅
4   for 每个锚点 a ∈ A do
5       G' ← G' ∪ 加权双向 BFS 扩展(a)     // 受深度/边权重/节点上限约束，
6   end for                                 // 候选按复合评分排序（式 3-4）
7   S ← 距离感知骨架抽取(G')               // 按到锚点距离分级保留信息
8   C ← 符号化序列化(S)
9   while tokens(C) > B do                 // Token 预算迭代控制
10      按五级降级链降级（降压缩档→减节点→减深度→提边权→仅锚点）
11      重新执行 7-8
12  end while
13  return C
\end{verbatim}

\textbf{（2）置信度衰减加权的锚点选择}。锚点（anchor）是上下文生成的起点节点。给定自然语言任务描述，系统通过语义检索在图谱中定位初始锚点候选；对于业务能力类任务，以能力节点（C 层）为入口，经 business\_map 边扩展至结构层锚点。能力节点 $c$ 的置信度 $\mathrm{Conf}_C$ 由其关联的 business\_map 边证据强度聚合而来：

\begin{equation}
\mathrm{Conf}_C = 1 - \prod_{e \in E_c}(1 - W_e) \tag{3-2}
\end{equation}

其中 $E_c$ 为能力节点 $c$ 关联的 business\_map 边集合，$W_e$ 为边 $e$ 的聚合权重。$\mathrm{Conf}_C$ 越高，表明该能力与代码结构的映射越被多源证据确认。

锚点扩展面临模块膨胀问题：若某能力的置信度已经很高，继续向 L1 模块层扩展锚点会引入大量低相关模块；反之若置信度低，仅靠能力入口召回不足。为此引入置信度衰减加权：

\begin{equation}
w_{L1} = \exp(-\alpha \cdot \mathrm{Conf}_C), \quad \alpha = 3.0 \tag{3-3}
\end{equation}

其中 $w_{L1}$ 为 L1 层锚点的扩展权重，$\mathrm{Conf}_C \in [0,1]$ 为能力置信度。高置信能力抑制模块层扩展（其映射已足够聚焦），低置信能力放开扩展以保证召回；冷启动场景（能力规范缺失、Conf\_C 不可得）时退化为纯结构检索，不阻塞流程。

\textbf{（3）加权双向 BFS 子图裁剪}。从多个锚点出发，在图上执行加权双向广度优先搜索：既沿出边（下游依赖）也沿入边（上游依赖）扩展，支持经 business\_map 边跨层扩展。选择 BFS 而非其他图遍历策略基于任务特性考量：其一，相对于 DFS--深度优先搜索容易沿单条依赖链进入局部深层的传递依赖（如工具函数的层层调用），偏离任务主体，而代码任务的影响范围通常集中在锚点附近的浅层邻域；其二，相对于 Dijkstra 类最短路径算法--其精确距离计算在大图上的成本较高，而上下文生成对"跳数级别的近似距离"已足够敏感（距离衰减因子按跳数设计）；其三，BFS 按层扩展的特性天然与"距离锚点越近越重要"的骨架分级假设（见（4））一致，且扩展深度即为预算控制参数，二者统一。

扩展过程中每个候选节点按复合评分排序：

\begin{equation}
\mathrm{score}(v) = \lambda_s \cdot \mathrm{sem}(v) + \lambda_t \cdot \mathrm{struct}(v), \quad \lambda_s = 0.6,\ \lambda_t = 0.4 \tag{3-4}
\end{equation}

\begin{equation}
\mathrm{sem}(v) = \mathrm{sim}(q, v) \cdot \exp(-\beta \cdot d(v)), \quad \mathrm{struct}(v) = \mathrm{deg}(v) \text{ 归一化}
\end{equation}

其中 $\mathrm{sim}(q,v)$ 为任务描述与节点向量的余弦相似度，$d(v)$ 为节点到最近锚点的跳数（距离衰减因子 $\beta$ 控制语义分随距离的衰减），$\mathrm{struct}(v)$ 为节点结构重要度（以度数归一化近似）。裁剪受三类预算约束：最大深度（默认 3）、最小边权重（默认 0.7，低于该权重的边不沿其扩展）与节点上限（默认 100）。

\textbf{（4）距离感知骨架抽取与符号化序列化}。裁剪后的子图仍包含大量与锚点距离较远的细节节点，直接序列化会造成信息稀释。骨架抽取按节点到锚点的距离分级保留信息：距离 0（锚点）保留完整属性与签名；距离 1 保留标准信息（名称、类型、关键属性）；距离 ≥2 仅保留最小骨架（名称与层级）。序列化输出采用层级符号化表示，以 ⊃ 表达包含、-> 表达依赖、⇄ 表达双向关联，◉ 标记锚点，使同等信息量下的 Token 占用显著低于自然语言描述或 JSON 全量输出。

在骨架分级基础上提供三档压缩等级（loose/standard/extreme），档位递进时提高距离阈值、削减保留属性、收紧 Token 上限，供上层按任务预算选择，三档的参数对比如表 3-4 所示；其压缩率、上下文召回率与信息保真度的权衡，以及三档对应的度量指标（Token 缩减率、上下文召回率、下游任务完成效果）在第 5 章消融实验（消融 B）中量化评估。

表 3-4 三档压缩策略参数对比

\begin{table}[htbp]
\centering
\small
\begin{tabularx}{\textwidth}{|X|X|X|X|}
\hline
\textbf{压缩档位} & \textbf{节点规模（默认）} & \textbf{信息保留粒度} & \textbf{典型 Token 上限（参考）} \\
\hline
loose & 约 100 节点 & 锚点与距离 1 节点保留较完整属性与签名 & 约 8K \\
standard & 约 60 节点 & 锚点完整、距离 1 标准、其余骨架 & 约 4K \\
extreme & 约 30 节点 & 仅锚点标准信息，其余最小骨架 & 约 2K \\
\hline
\end{tabularx}
\end{table}

注：表中节点数与 Token 上限为默认参考值，实际由 Token 预算迭代控制（见（5））动态决定；具体数值以第 5 章实验实测为准。

\textbf{（5）Token 预算迭代控制}。给定用户或上层指定的 Token 预算，序列化输出若超预算，按五级降级链迭代执行直至满足：降低压缩档位 -> 削减节点数 -> 削减扩展深度 -> 提高边权重阈值 -> 仅保留锚点；锚点在任何降级层级下始终保留，保证上下文不失去任务焦点。此外流水线设三级系统级降级：向量索引缺失时退化为结构检索；图谱缺失时退化为文件系统提示；序列化失败时输出锚点清单，确保上下文生成环节不阻塞开发流程。


\section{本章小结}
本章阐述了 wpw 的知识基础设施构建方法。在知识库层面，以能力规范为核心组织稳态知识，通过归档沉淀机制实现过程知识向稳态知识的转化，形成可持续积累的项目级知识库。在知识图谱层面，提出能力-代码双层四层本体模型，其中业务能力节点作为上层语义锚点与代码结构层通过 business\_map 边映射而非包含；业务-代码关联由文档提取、命名匹配、语义匹配与 Git 历史四类证据源经 noisy-OR 概率融合自动建立，每条关联可溯源至最权威证据，并以权重剪枝控制噪声。在大模型结合层面，提出区别于文本检索 RAG 的结构感知上下文生成流水线：置信度衰减加权锚点选择、加权双向 BFS 子图裁剪、距离感知骨架抽取与三档压缩序列化，配合五级降级的 Token 预算迭代控制，在预算约束下生成任务相关、结构连贯、可解释的高密度代码上下文。本章方法的有效性将在第 5 章通过上下文 Token 消耗与压缩率、上下文召回率与下游任务完成效果，以及图谱贡献消融实验（消融 A/B 与组件级基线）予以验证。

\textbf{本章参考文献}（沿用第 1 章编号，全文合稿时统一）：

\begin{itemize}
\item {}[4] Liu N F, et al. Lost in the Middle. arXiv:2307.03172, 2023.
\item {}[5] Lewis P, et al. Retrieval-Augmented Generation for Knowledge-Intensive NLP Tasks. NeurIPS 2020.
\item {}[11] Guo D, et al. GraphCodeBERT: Pre-training Code Representations with Data Flow. ICLR 2021.
\item {}[13] Cleland-Huang J, et al. A Heterogeneous Solution for Improving the Return on Investment of Requirements Traceability. IEEE RE 2004.
\item {}[14] Borg M, et al. Recovering from Requirements to Code: An Empirical Study of Traceability Link Recovery. EMSE 2014.
\item {}[15] Sultanov H, Hayes J H. Application of Genetic Algorithms to the Requirements Traceability Problem. REJ 2010.
\item {}[18] Yamaguchi F, et al. Modeling and Discovering Vulnerabilities with Code Property Graphs. IEEE S\&P 2014.
\end{itemize}
\textbf{本章图表清单}（成稿时绘制）：

\begin{itemize}
\item 图 3-1 能力-结构双层图谱本体模型（C / L1 / L2 / L3 与四类边）
\item 图 3-2 多源证据融合与上下文生成流水线（证据收集 -> noisy-OR 聚合 -> 锚点 -> BFS -> 骨架 -> 序列化）
\item 图 3-3 wpw 知识库整体架构（数据来源层 -> 知识组织层 -> 智能体消费层）
\item 图 3-4 business\_map 生成流程（能力节点 -> 候选召回 -> 四类证据 -> 融合 -> 过滤 -> 边）
\item 算法 3-1 图谱驱动的任务上下文生成（伪代码，已含）
\item 表 3-1 多源证据体系（已含）
\item 表 3-2 图谱边类型（已含）
\item 表 3-3 L3 代码元素节点属性定义（已含）
\item 表 3-4 三档压缩策略参数对比（已含，具体数值以第 5 章实验实测为准）
\end{itemize}
